\section{Conclusion}\label{chap4}
In this last chapter, we provide a brief summary of major topics convered in 
the thesis and assess the impending future work. 

\subsection{Summary}
In this thesis we,
\begin{itemize}
    \item studied the physical phenomenon of flash evaporation. We reviewed 
    the contradictory effects of low-pressure and low-gravity environments to liquid-gas
    phase change phenomena. We review the existing flash evaporation models and their 
    capabilities, limits, and applicabilities to different flow scenarios.
    %The increased superheat limit due to increased
    %delay in nucleation is counteracted by the higher rate of heterogeneous
    %nucleation on surfaces. Bubble shapes tend to more ideal descriptions, and the lack
    %of buoyancy or strong convection in the system eliminates complex bubble interactions.
    %We described the limiting rate of phase-change using bubble-growth and kinetic theory.
    Based on experimental 
    observations and review of literature, we identified the governing 
    factors in an accurate modelling of flash evaporation.

    \item revisted the multiphase flow models for two-phase liquid-gas systems and 
    detailed the context of jump conditions in the models. We also discuss turbulence and
    describe the numerical implementation of these models in STAR-CCM+.
    
    \item assessed the complex maze of closure relations for the interphase interactions
    in the two-fluid model and arrived at a set of generally applicable closures based on
    physical principles instead of empirical coefficients. 

    \item modelled the 1D Stefan problem and the 2D Rayleigh-Taylor instability. This was done
    to assess the treatment of a moving 
    interface in the context of a fixed grid method like the finite-volume method. The 
    Stefan problem also serves as a benchmark to test for the accurate modelling of the phase change phenomenon.
    The Rayleigh-Taylor instability is a fundamental phenomenon in two-phase flows with a significant 
    density difference. The limitations of our assumptions in context of high Reynolds numbers
    and turbulence are discussed in 
    this context.
    
    \item described the mathematical formulation of the infiltration of water in a low-pressure
    chamber which was subsequently modelled using the Volume-of-Fluid method.
    We implemented adaptive mesh refinement based on user-defined field functions in order to
    refine the mesh locally during flow evolution. We finally analysed the flow evolution and 
    determined the initial conditions expected in the flash evaporation subproblem.    
\end{itemize}

\subsection{Outlook}
While this thesis investigated the factors affecting a modelling framework for flash evaporation, the 
 problem remains unmodelled and forms the immediate next steps. 
 \begin{itemize}
    \item An assessment of the effects of RANS turbulence models is required for flash evaporation applications.
    Investigations combining a hybrid RANS-LES approaches could be suitable possibilities.
    \item Building on the validated treatment 
 of two-phase flows using the VOF method, investigations could be performed identifying 
 the possibility of modelling phase change using phase specific source terms in STAR-CCM+. Equilibrium
 conditions can be enforced by disregarding the energy equation solver and implementing an additional 
 scalar equation to treat temperature. 
    \item An ideal flash evaporation model would account for the homogeneous and heterogeneous nucleation rates
    along with phase transfer at the interface governed by the kinetic theory. To this end, a scalar equation 
    solving for the conservation of the number of bubbles in the system can be solved. This would enable
    consistent interfacial area concentration and closure feedback for heat, mass, and momentum exchange.
 \end{itemize}

\subsubsection{A Theoretical Stipulation}
To model flashing under rapid depressurisation with thermal and velocity non-equilibrium, 
we propose an Eulerian Multiphase (two-fluid) formulation 
with Large-Scale Interface capabilities, coupled to interfacial 
area concentration (IAC) transport, a bubble population balance for 
number density and nucleation, and a kinetic interfacial mass-flux closure of Schrage type.

\paragraph{Governing phasic balances and sources.}
For each phase \(i\in\{l,v\}\), solve the volume-averaged mass, 
momentum, and energy balances as in the two-fluid model with a mass source in the governing equations to the 
the kinetic mass transfer using the Schrage equation. 

\paragraph{Interfacial area concentration transport.}
Evolve the IAC field \(a_{lv}\) to reflect transient nucleation, growth, coalescence/breakage, and phase change:%
\[
\frac{\partial a_{lv}}{\partial t} + \nabla\cdot(a_{lv}\,\boldsymbol{u}_{mix})
\;=\; S_{nuc} + S_{grw} - S_{coal} - S_{brk} \;+\; S_{pc},%
\]
where \(S_{nuc}\) adds area from newly nucleated bubbles, \(S_{grw}\) accounts for radius change \( \dot{R}\), \(S_{coal}\) and \(S_{brk}\) model topology changes, and \(S_{pc}\) captures interfacial recession/advance due to phase change at large interfaces (LSI regime).%
In dispersed-bubble cells, relate \(a_{lv}\) to the bubble size distribution \(n(d)\) by%
\[
a_{lv} \;=\; \int_0^\infty \pi d^2\,n(d)\,\mathrm{d}d
\quad\text{(spherical bubbles)}.%
\]

\paragraph{Population balance for bubble number density.}
Solve a number-density transport for diameter \(d\) (or radius \(R\)) with growth velocity \(G=\dot{d}\), 
coalescence/breakage kernels, and nucleation sources:%
\[
\frac{\partial n}{\partial t} + \nabla\cdot(n\,\boldsymbol{u}_v) + \frac{\partial}{\partial d}\!\bigl(G\,n\bigr)
\;=\; B_{nuc}(d) + B_{coal}(d) + B_{brk}(d) \;-\; D_{coal}(d) - D_{brk}(d).%
\]


\paragraph{Bubble growth law and thermal coupling.}
Use a composite growth rate that superposes inertial- and thermal-limits, consistent with diffusion-controlled Stefan scaling at later times:%
\[
\dot{R} \;=\; 
\sqrt{\frac{2\,(p_{sat}(T_l^{int})-p_l)}{3\rho_l}}
\;+\; 
\frac{\rho_l\,c_{p,l}\,|T_l^{int}-T_{sat}(p_l)|}{\rho_v\,h_{lv}}
\sqrt{\frac{3\,\alpha_l}{\pi\,t}}\,,
\]
and couple to \(\dot{m}\) by enforcing local interfacial energy balance and consistent choices of \(T_l^{int}\), \(T_{int}\), and heat-transfer coefficients for \(q_i^{int}\).%


\paragraph{Algorithmic outline for the flashing stage.}
\begin{enumerate}
  \item Initialise from the infiltration end-state: fields \((\alpha_i,\boldsymbol{v}_i,T_i,p)\), set vacuum/outflow boundary and depressurisation profile.%
  \item Enable EMP with multiregime blending and activate PBE+IAC in dispersed regions; seed initial \(n(d)\) from known particulate load or a uniform low background consistent with silicate content.%
  \item Compute nucleation source \(B_{nuc}\) using CNT with heterogeneous correction where walls/impurities dominate; map to \(n(d)\) and update \(a_{lv}\).%
  \item Advance \(\dot{R}\) and \(G=\dot{d}\) with thermal coupling; update \(a_{lv}\) and \(n(d)\) including coalescence/breakage.%
  \item Evaluate kinetic \(\dot{m}\) (Schrage) and apply \(\Gamma_i=a_{lv}\dot{m}\), \(\dot{Q}\), and interfacial momentum sources in phasic balances; iterate to convergence per time-step.%
  \item Remesh/refine adaptively near rapid void-fraction transients and LSI interfaces; limit interface Courant number and apply interface-turbulence damping if needed.%
\end{enumerate}

This plan provides a theoretically working 
NHNEM pathway with physics-based interfacial exchanges, retaining slip 
and non-equilibrium, while remaining compatible with STAR-CCM+ EMP/LSI numerics.%
